\section{Application Case Study: AI-Inference
Serving}

To show the framework instantiates a contemporary problem, we cast
request routing in an AI-inference serving system as online capacitated
matching: resources are model replicas / cache shards serving up to
\(c\) concurrent requests, arrivals are requests drawn from a
(non-uniform, bursty) traffic distribution, and an edge is a capability
or cache affinity. Two quantities must be kept apart: raw \emph{goodput}
is the fraction of arrivals served, while the \emph{competitive ratio}
is the number served divided by the capacitated optimum on the same
realized instance; they coincide only when the optimum can serve every
arrival, which under overload it cannot, so we report the competitive
ratio. We instantiate the mapping across four serving concerns: capacity
\(c\) (on a synthetic topology, where advice quality can be swept
directly), stale traffic forecasts (Wikipedia pageviews), dynamic
service times (an Azure LLM inference trace), and prefix-cache-aware
routing (the Mooncake prefix-cache trace).

We present this as a \textbf{case study, not a novelty claim} --- the
systems facts we recover are established, and the ``with-predictions''
vocabulary is a re-labeling of them. The framework does reproduce them
cleanly: capacity is a \emph{substitute} for algorithmic robustness
(blind trust in a forecast degrades further as capacity grows, while a
capacity-aware baseline stays safe); a live-load signal beats a stale
forecast under dynamic service times; and a \emph{stable} cache-affinity
router beats a reactive one --- the reverse of the traffic-forecast
case. Because the literature (§9) suggested the with-predictions lens
might yield a \emph{new} actionable serving result on a tail objective,
we probed it: on an SLO/tail objective under bursty load, a
non-predictive policy (static headroom or reactive-adaptive reservation)
comes within \(\le 3\%\) of a policy with perfect foresight across every
regime we swept. We found no regime where foresight helps --- the tail
objective is forgiving too, a third face of the paper's wall. Serving
therefore stays a case study.

\section{Related Work}

Algorithm attributions are given where each algorithm is introduced
(§2), and the positioning of our budget--stakes law against prior
testing and tradeoff results is in §7; here we place the remaining
landscape.

\textbf{Scope of our novelty claims.} Both claims we mark as new --- the
budget--stakes law, and the observation that testing the \emph{payoff}
separates from testing the prediction's distance --- rest on a targeted
search rather than an exhaustive one, and we state its extent so it can
be checked. We searched the learning-augmented /
algorithms-with-predictions literature through \textbf{April 2026}, over
DBLP, arXiv (cs.DS, cs.LG) and Google Scholar, for prefix- or
sample-based acceptance rules --- keywords \emph{test-before-trust},
\emph{advice testing}, \emph{trust the prediction},
\emph{test-then-commit}, \emph{prediction validation},
\emph{consistency--robustness trade-off}, \emph{tolerant testing} ---
and read in depth the four nearest lines (test-before-trust; switching
and combining frameworks; distributional advice of unknown quality;
data-driven algorithm selection), including the test-before-trust
group's own April-2026 survey talk of their programme. We found no
acceptance rule that estimates the payoff of following. A citation trace
of every work citing \cite{choo2024imperfect} would be the completionist
step and we have not performed it; the claims should be read as scoped
to this search.

\textbf{Learning-augmented algorithms.} The consistency/robustness
framework originates with competitive caching with predictions
\cite{lykouris2018caching} and the optimal-tradeoff analyses that
followed \cite{weizhang2020tradeoffs}; our thesis --- that on
average-case matching the value is robustness, not consistency --- is a
same-spirit but problem-specific statement, made quantitative and then
proven necessary. The direct experimental-study template is Chłędowski
et al.~\cite{chledowski2021caching} in caching, whose
blind-follow-with-switching combiner we benchmark (§5.5); the switching
genre extends to prediction-quality-triggered frameworks such as
SemiTrust-and-Switch \cite{antoniadis2025switching}, and distributional
advice of unknown quality is handled by hedging in ski rental
\cite{cui2026skirental} and in Diakonikolas et al.'s sampling-access
model --- among the works covered by the search above, none tests the
payoff and none commits once. The most closely related active line is
the \emph{test-before-trust} program of Choo, Gouleakis and coauthors
\cite{choo2024imperfect,bhattacharyya2025product}, which statistically
validates advice before use --- always by distributional distance; our
§5.4/§7.7 argue the decision-relevant statistic is the payoff instead.

\textbf{Algorithm selection, best-arm identification, and sequential
testing.} Choosing between two policies from samples at gap \(g\) costs
\(\Theta(1/g^2)\) observations --- the classical engine of sequential
analysis \cite{wald1947sequential}, of data-driven algorithm selection
\cite{gupta2017pac,balcan2020datadriven}, of two-armed best-arm
identification \cite{mannortsitsiklis2004,kaufmann2016complexity}, and
of off-policy evaluation \cite{dudik2011doubly}. We import that engine
and claim no novelty for it. The delta is the \emph{observation model}:
in algorithm selection each sample is a whole instance with directly
observable performance; a bandit pull directly reveals the pulled arm's
reward; off-policy evaluation consumes logged actions with known
propensities. Our prefix is none of these --- it is a \textbf{passive
stream of arrivals}, containing no actions and no rewards of either
policy. That such a stream reveals a policy pair's full-horizon value
gap at all, at rate \(\sigma^2/g^2\) and --- by the law's lower half ---
provably no faster under \emph{any} rule, is the content of the payoff
identity and of Theorem 1, not of the statistics; and the bootstrap
calibration of §5.4 is what survives of the identity on benchmark
instances with capacity kinks.

\textbf{Online matching with predictions.} MinPredictedDegree
\cite{aci2022mpd} and the test-and-fallback schemes
\cite{choo2024imperfect,bem2026testmatch} are the algorithms we unify
and, for the latter, bound; each was previously studied in isolation.
Blending designs for the fractional relaxation \cite{choo2025fractional}
optimize the robustness--consistency Pareto frontier directly and test
nothing --- complementary to, and untouched by, the budget--stakes law,
which prices the \emph{decision to test}. Borodin et
al.~\cite{borodin2018experimental} is the advice-free experimental
foundation our benchmark extends.

\textbf{Distribution testing.} Tolerant identity testing
\cite{canonne2022tolerance,valiant2011unseen} and \(\ell_1\)-distance
estimation \cite{jiao2018l1} enter our story as the explanation of why
the field's empirical-\(\ell_1\) acceptance rules are blind at sublinear
prefixes --- the near-linear \(\tilde\Theta(r/\log r)\) price of
tolerance is real. Our budget--stakes law shows the follow/fallback
\emph{decision} nonetheless escapes that barrier: its decision-relevant
statistic is the payoff, not the distance (§7.7), which is why our lower
bound is a Hellinger computation rather than a testing reduction.

\textbf{Serving systems.} The systems results our §8 case study recovers
are established across Preble, Mooncake, SageServe, and related work; we
use them as ground truth, not as contributions.

\section{Limitations and Conclusion}

\textbf{Limitations.} (i) The experiments are in the known-i.i.d. model.
Every known-i.i.d. instance is also a random-order instance, so
guarantees proved in the random-order model carry over to ours but not
conversely; the law itself does extend to random order (Remark, §7.8),
but our empirical wall is an average-case statement and we do not claim
it for adversarial arrival order. (ii) Following the original authors,
the test-and-fallback experiments use an empirical-\(\ell_1\) surrogate
for the (unimplemented) distribution tester; the theory does \emph{not}
inherit this modeling --- the law binds any rule --- and §7.7 separately
explains the surrogate's blindness. (iii) The degree- and
histogram-prediction families do not map onto every graph, which is why
we report them in parallel panels rather than one table. (iv) Each real
modality is exercised by one trace. (v) The budget--stakes law is proven
on decomposable (disjoint-cell) families, where the payoff identity
holds; whether a \emph{non}-decomposable family can push the budget
above \(1/\delta^2\) is open. Our novelty claim for payoff-estimating
acceptance rules rests on the search reported in §9: that search is
complete, but bounded by the venues, keywords and cut-off stated there,
so we present the claim as scoped rather than exhaustive. An earlier
draft claimed a tolerant-testing impossibility for any sublinear rule;
that claim was refuted during its own witness step and the refutation is
reported in §7.7.

\textbf{Conclusion.} We gave the first unified experimental study of
learning-augmented online bipartite matching --- advice-free baselines,
MinPredictedDegree and its augmentations, and the test-and-fallback
algorithms --- on one harness, with confidence intervals, across
synthetic graphs, six real graphs, and real traces. Across every setting
the same wall appears: the advice-free baseline is already near-optimal,
unguarded prediction-following crashes below it, and the value of the
sophisticated algorithms is downside protection, delivered by structural
or adaptive robustness with distinct trade-offs. We then priced the wall
for the test-and-fallback class with a two-sided budget--stakes law: the
follow/fallback decision costs a prefix of
\(\tilde\Theta(\sigma^2/g^2)\) --- baseline slack over squared stakes
--- for \emph{any} decision rule, an explicit directional test achieves
it (implemented and benchmarked, it avoids both threshold pathologies,
§5.4), and --- because stakes are capped by the baseline slack --- on
strong-baseline instances the price exceeds the horizon: upsides below
\(\Theta(\sqrt{(1-\rho_{\mathrm{base}})/n})\) are uncapturable at any
feasible prefix. Experiments discover the wall; theory sends the bill.
The practical message is a single sentence --- \textbf{on average-case
matching, the advice's upside is smaller than the price of finding out
whether to trust it} --- with one constructive corollary for algorithm
designers: \emph{test the payoff, not the prediction} (the payoff is a
per-sample observable where the distance is not). Progress on
predictions for online matching will come from the regimes this paper
brackets: adversarial or non-stationary arrivals, and objectives (tail,
fairness, cost) where the advice-free baseline is genuinely far from
optimal.
